\begin{claim}{C.75}{Proved}{The actual orbit strips for a one-dimensional proper operator action}
Let $O$ be a plane, including an actual open disk, and
let $Q\leq\operatorname{Homeo}(O)$ be closed, locally
compact, free and proper, with $\dim Q=1$.
Let $f_\tau$ be a continuous based loop in its
centralizer, and put $G=\langle f_\tau\rangle$.
Retain the actual principal coordinate of C.33,
with its inverse, in which
\[
 \begin{split}
 O&\cong\mathbb R_s\times\mathbb R_t,\\
 f_\tau(s,t)&=(s+u_\tau(t),\beta_\tau(t)),\\
 A&=\{t:\beta_\tau(t)=t\text{ for every }\tau\}.
 \end{split}                                                   \tag{C.75.1}
\]
The set $A$ is closed. On every component $I$ of
$\mathbb R\setminus A$, precisely one of the
following two descriptions holds.
Either $\mathbb R\times I$ is one open $G$-orbit,
or there is a continuous function $b_I:I\to\mathbb R$
such that all the orbits in that strip are exactly
\[
                   \{(c+b_I(t),t):t\in I\},
                                  \qquad c\in\mathbb R.        \tag{C.75.2}
\]
In the latter case the actual coordinate
$(s,t)\mapsto(s-b_I(t),t)$, with inverse
$(w,t)\mapsto(w+b_I(t),t)$, changes every restricted
generator to $(w,t)\mapsto(w,\beta_\tau(t))$.
The effective restricted action and its base image
are topologically isomorphic in their actual
compact-open and inverse topologies.

At each $a\in A$ the orbits on $\mathbb R\times\{a\}$
are individual points when every $u_\tau(a)$ vanishes,
and otherwise that entire horizontal line is one
orbit. This lists all orbits. In particular every
finite interior frontier line of an open orbit is
the actual proper line at an endpoint of its
base interval. There are no circular orbits.
The graph in (C.75.2) is proper in its indicated
strip; extension of $b_I$ to an endpoint is
not asserted.

There is a further exact comparison for a finite
endpoint $a\in A$ of a base interval $I$.
Write $g(s,t)=(s+u_g(t),\beta_g(t))$ for the
original element $g\in G$, and put
\[
 c_a(g)=u_g(a),\qquad
 G_t=\{g\in G:\beta_g(t)=t\}\quad(t\in I).
                                                               \tag{C.75.3}
\]
The character $c_a$ is continuous. There are
exactly three alternatives: $c_a=0$; or
$c_a\ne0$ and $c_a(G_t)=0$; or
$c_a(G_t)=\mathbb R$. In the second alternative
there is an actual homeomorphism
$b:I\to\mathbb R$ satisfying
\[
                 b(\beta_g(t))=b(t)+c_a(g)
                                  \qquad(g\in G,\ t\in I).
                                                               \tag{C.75.4}
\]
In the third alternative every real endpoint
displacement is realized by an actual source
element fixing the chosen base point $t$.
All component-operator equations and the kernel
of the comparison action used below are retained.
No continuous choice of these correcting source
elements, or frontier filling, is inferred.
\end{claim}

\begin{proof}
\emph{The principal coordinate and the original group words.}
C.33 gives a Lie group with identity component
$Q^0$, and its actual product principal bundle.
Because $O$ is a plane, neither factor in that
product can be a circle: a circle factor gives
a nontrivial fundamental-group factor, whereas
the plane has trivial fundamental group.
Thus $Q^0=\mathbb R$, its base is $\mathbb R$,
and an actual section gives the coordinate
and inverse in (C.75.1).
Every labelled element of $Q^0$ acts by
$(s,t)\mapsto(s+c,t)$.

Commutation with these translations forces
every $g\in G$ to have the stated form.
The base homeomorphism is increasing because
each generator is joined to the identity.
The inverse map is exactly
\[
 (s,t)\longmapsto
       \bigl(s-u_g(\beta_g^{-1}(t)),\beta_g^{-1}(t)\bigr).
                                                               \tag{C.75.5}
\]
For the composition convention $gh=g\circ h$,
\[
 \beta_{gh}=\beta_g\beta_h,\qquad
 u_{gh}(t)=u_h(t)+u_g(\beta_h(t)).
                                                               \tag{C.75.6}
\]
The functions and their inverse maps are jointly
continuous on every compact parameter family.

Put $P=\{f_\tau\}\cup\{f_\tau^{-1}\}$, retaining
the actual maps. It is compact, path connected,
and contains the identity. Every $P^n$ is compact
and path connected, and $\bigcup_nP^n=G$.
More explicitly an element
$g=f_{\tau_1}^{\epsilon_1}\cdots
f_{\tau_m}^{\epsilon_m}$ has the actual path
\[
 r\longmapsto
 f_{r\tau_1}^{\epsilon_1}\cdots
 f_{r\tau_m}^{\epsilon_m},\qquad 0\leq r\leq1,
                                                               \tag{C.75.7}
\]
from the identity to $g$, entirely in $P^m$.
No path in a larger closure is substituted.

The fixed set of all base generators is
exactly $A$, and it is also the fixed set
of their generated base group.
The base orbit of any $t\notin A$ is an
open interval. Indeed choose a based path
of base maps moving that point. That
path and its inverse give points on both
sides of it; their evaluation intervals
give a neighborhood in the orbit.
The same argument at each point of the
orbit makes the orbit open. It is an
interval because it is the increasing
union of the connected compact evaluation
images of $P^n$.
A finite endpoint of this orbit is fixed
by every increasing base map preserving
the orbit. Hence the orbit is exactly
the component $I$ of $\mathbb R\setminus A$
containing $t$.

\emph{The precise interval consequence of the word triangle.}
For a continuous path $h_r$ of actual
maps of the form in (C.75.1), starting at
the identity, suppose that $h_1(s,t)$
has second coordinate $t$.
The entire horizontal interval joining
$(s,t)$ to $h_1(s,t)$ is contained in
the orbit of $(s,t)$ under the group
generated by that path.
This is the exact C.62 word-triangle
argument with second-coordinate quotient.
Its hypotheses here can be checked directly:
the source is the actual product plane,
the real fibers are proper horizontal
lines, the induced base maps are
increasing, and the endpoint equality
is real equality.
No original-open-orbit hypothesis
enters that triangle proof.

More explicitly, subdivide the path by
the closed set of parameters at which
its second-coordinate evaluation equals
$t$. On an excursion $(a,b)$, put
$A_0=h_a(s,t)$ and $k_r=h_rh_a^{-1}$.
The paths $k_rA_0$ and
$k_bk_r^{-1}A_0$ have opposite strict
second-coordinate signs relative to
the original fiber. Their concatenation
has winding $1$ or $-1$ about every
interior point of the interval between
their two endpoints, as in C.62.9.
The triangle $(v,u)\mapsto k_uk_v^{-1}A_0$
therefore contains that interval.
Interpolating on the closed zero set
and these excursion gaps, with uniform
continuity on their compact parameter
set, gives the complete endpoint
interval, exactly as in C.62.
Every displayed value is an actual
three-word value from the path.
For (C.75.7) it consequently belongs
to $P^{3m}(s,t)$.

\emph{The two possible strip actions.}
For each $t\in I$, define the actual
additive subgroup
\[
       D_t=\{u_g(t):g\in G_t\}\leq\mathbb R.
                                                               \tag{C.75.8}
\]
Additivity and inverses follow from
(C.75.5--6). If $h$ takes the base point
$t$ to $t'$, conjugation by that actual
map takes a shift $c\in D_t$ to the
same shift $c$ in $D_{t'}$:
its first-coordinate translations
cancel. Thus $D_t=D_{t'}$ throughout
this base orbit.

If $D_t$ contains a nonzero $c$, take
its actual source element $g$ and the
path (C.75.7). The preceding triangle
puts every point of the interval from
$(0,t)$ to $(c,t)$ in the original orbit.
Each such orbit element fixes the
base point, so $D_t$ contains the
full interval from zero to $c$.
An additive subgroup containing a
nondegenerate interval is all of
$\mathbb R$. Transitivity on the
base interval, followed by these
arbitrary fiber translations, proves
that $\mathbb R\times I$ is one orbit.

Suppose instead that $D_t=0$.
Choose $t_0\in I$ and $x_0=(0,t_0)$.
Every fiber over $I$ meets $Gx_0$
exactly once: existence follows
from base transitivity, and two
such points would give a nonzero
element of $D_t$ after composing
their source words.
Write that point as $(b_I(t),t)$.

This function is continuous with
its actual, not a postulated,
graph topology. The compact set
$P^n(x_0)$ projects continuously
and bijectively to a compact
interval $I_n$. A compact-to-Hausdorff
bijection is a homeomorphism.
The intervals $I_n$ increase and
exhaust $I$. For each $t\in I$,
choose $t_-<t<t_+$ in $I$.
Both belong to some $I_n$, so
$[t_-,t_+]\subset I_n$.
Thus the inverse projection,
and hence $b_I$, is continuous
on a neighborhood of every $t$.

Every horizontal translation
commutes with $G$. Its translates
of $Gx_0$ are all the other
orbits, proving (C.75.2).
Invariance of this actual graph
under $g$ gives
\[
           u_g(t)=b_I(\beta_g(t))-b_I(t).
                                                               \tag{C.75.9}
\]
This verifies both conjugacy
directions of the asserted
coordinate change.
The graph is closed in
$\mathbb R\times I$, and its
projection is a homeomorphism
onto $I$, so it is proper
in that strip.

Under this coordinate, the
effective restricted group
is exactly $\{\,\id_{\mathbb R}
\times\beta:\beta\in\beta(G)|_I\,\}$.
Projection to the base has
the explicit continuous inverse
$\beta\mapsto\id\times\beta$,
also on inverse maps, in the
compact-open topologies.
Evaluation on the section
$\{0\}\times I$ proves
continuity in the other direction.
This proves the stated topology
of the effective images. It
makes no quotient-topology
claim for restriction from
the original global group.

For $a\in A$, the first-coordinate
shifts form the image of the
path-connected group $G$ under
the continuous evaluation
$g\mapsto u_g(a)$.
This is a connected additive
subgroup of $\mathbb R$, hence
zero or $\mathbb R$.
The zero alternative is exactly
the vanishing of every generator
$u_\tau(a)$. This proves the
remaining orbit description.
The frontier inside the product
plane of an open strip is
exactly the horizontal lines
at its finite endpoints.

\emph{The endpoint character with all source labels retained.}
Fix the indicated finite
endpoint $a$ and interval $I$.
By (C.75.6), $c_a$ is an
additive character, continuous
in the original group topology.
Form the labelled comparison
action on a second product:
\[
 \begin{split}
 T_a:G&\longrightarrow
                 \operatorname{Homeo}(\mathbb R_z\times I),\\
 T_a(g)(z,t)&=(z+c_a(g),\beta_g(t)),\\
 \ker T_a&=\{g:c_a(g)=0,\ \beta_g|_I=\id\}.
 \end{split}                                                   \tag{C.75.10}
\]
The displayed kernel is kept;
no faithfulness or closedness
of this image is assumed.
Every finite word and every
path (C.75.7) has precisely
its displayed image under $T_a$.

The strip proof just given
therefore applies to this
actual image action. Its
fiber stabilizer at $t$ is
exactly $c_a(G_t)$, so it is
zero or all of $\mathbb R$.
In the latter case the same
triangle supplies actual
source words with every
intermediate real displacement,
and their finite products
give every real displacement.
They fix this specified base
point; a continuous choice
of these words has not been
constructed.

In the zero case the actual
graph function from that
argument satisfies
\[
           b(\beta_g(t))=b(t)+c_a(g).
                                                               \tag{C.75.11}
\]
Suppose that $c_a$ is nonzero.
Its path-connected additive
image is then $\mathbb R$,
so $b(I)=\mathbb R$ by base
transitivity and (C.75.11).
We check that $b$ is injective.
If two distinct arguments
had the same value, either
$b$ would be constant on
the interval between them,
or it would have an interior
local maximum or minimum
there.
An open interval of constancy
is transported by the
transitive base action to
neighborhoods of every point,
which would make $b$ constant.
A local maximum is similarly
transported to every point.
A continuous real function
having a local maximum at
every point has countable
image: each of its values
is the maximum on one
rational closed interval
inside a local maximum
neighborhood. Its connected
image must therefore be
a singleton. The identical
argument with minima handles
the other case.
All alternatives contradict
$b(I)=\mathbb R$.
Thus $b$ is a continuous
injective surjection from
an interval to $\mathbb R$,
hence a strictly monotone
homeomorphism. This proves
(C.75.4) with its exact
inverse coordinate.

\emph{Component operators and the remaining boundary step.}
C.33 retains either no
nontrivial component operator
or one labelled generator
\[
       q(s,t)=(\lambda s+\kappa,t+d),
                    \qquad d>0,\quad\lambda\ne0 .
                                                               \tag{C.75.12}
\]
The original commutation
equations, with the constant
$\kappa$ retained, are
\[
 \beta_g(t+d)=\beta_g(t)+d,\qquad
 u_g(t+d)=\lambda u_g(t).
                                                               \tag{C.75.13}
\]
They give $A+d=A$, carry
each strip and its orbit
description to the translated
strip, and send its endpoint
character to $\lambda c_a$.
The actual maps $q$ and
$q^{-1}$ prove both transport
directions. Negative $\lambda$
is included.

All finite interior frontier
strata in this reduction
are consequently the named
horizontal lines, with the
original translation
characters just computed.
An actual disk coordinate
can still have complicated
limits at its original
outer circle, as C.74
demonstrates. Neither the
graph coordinate nor the
endpoint correction
surjectivity supplies a
jointly continuous filling
at that circle or across
the nonfixed interior
frontier lines. Those
extension steps remain
part of C.12.
\end{proof}
